\chapter{Implementación de un instrumento virtual básico}
\section{\obj}
\capacidad
\begin{itemize}
\item Implementar un instrumento virtual básico.
\item Implementar una interfaz de usuario básica.
\item Utilizar un DAQ para medir corriente y voltaje.
\end{itemize}

\section{\mat}
\textbf{A suministrar por la Escuela:}
\begin{itemize}
\item 1 dispositivo de adquisición de datos NI myDAQ

\end{itemize}
\textbf{A suministrar por el estudiante:}
\begin{itemize}
\item 1 breadboard
\item 1 potenciometro de \SI{10}{\kilo\ohm}
\item 1 resistor de \SI{220}{\ohm} y \SI{0.25}{\watt}
\item Cables de interconexión (jumpers) macho-macho
\end{itemize}


\section{\pro}
\begin{enumerate}
\item Realice las conexiones del circuito tal y como se indica en la Figura \ref{fig:L1F1}.
\begin{figure}[H]
    \centering
    \begin{circuitikz} 
        \draw
        (5.5,3) node[ground]{}
            to[V] 
        (5.5,6) node[above]{$+$\SI{15}{\volt}} 
            --
        (3,6)
        	to[american potentiometer,n=mypot, l_=0$\sim$\SI{10}{\kilo\ohm},i=$i$]
        (3,3) 
        %(mypot.wiper) to[short,-o] ++(2,0)
        (3,3)
            to[R, l_=\SI{200}{\ohm}, v^=$v$]
        (3,0)
            --
        (3,-0.5)
            --
        (5.5,-0.5) node[ground]{}node[above]{AGND}
        ;
        \draw 
        (5.5,1.5) node[op amp, noinv input up](opamp){AI0}
        (opamp.+) --++ (-0.5,0) --++ (0,0.7) to[short,-*]++ (-0.8,0)
        (opamp.-) --++ (-0.5,0) --++ (0,-0.7) to[short,-*]++ (-0.8,0)
        ;
        \draw[dashed,blue]
        (4,8) -- (7,8)node[midway, below]{myDAQ} -- (7,-1.5) -- (4,-1.5) -- (4,8);
    \end{circuitikz}
    \caption{Medición de corriente y voltaje en un circuito.}
    \label{fig:L1F1}
\end{figure}
\item Cree un instrumento virtual (\emph{.vi}) que realice lo siguiente:
    \begin{itemize}
        \item Leer en un While 100 muestras del valor de AI0 durante \SI{100}{\milli\second} y las promedie
        \item Calcular el valor promedio de las muestras obtenidas $v$ cada \SI{100}{\milli\second} y lo muestre en un gráfica en el panel frontal
        \item Calcular el valor promedio de las muestras obtenidas $i$ cada \SI{100}{\milli\second} y lo muestre en un gráfica en el panel frontal
        \item Iniciar la gráfica vacía. 
        \item Detener la toma de datos después de \SI{10}{\second}
    \end{itemize}
\item Una vez realizado el instrumento virtual córralo, modifique de forma aleatoria el valor del potenciometro y guarde los datos de las gráficas como archivos \emph{.csv}.
\item Incluya en su informe lo siguiente:
    \begin{itemize}
        \item Una captura de pantalla del diagrama de bloques
        \item Una captura de pantalla del panel frontal
        \item Una gráfica de $i$ realizada a partir del archivo \emph{.csv}, usando Python. Incluya los títulos de los ejes de forma adecuada. 
        \item Una gráfica de $v$ realizada a partir del archivo \emph{.csv}, usando Python. Incluya los títulos de los ejes de forma adecuada. 
    \end{itemize}
\end{enumerate}